% Standalone problem document, generated from the adjacent README.md.
% Compile with XeLaTeX. Mathematical content is maintained in Markdown.
\documentclass[11pt,a4paper]{article}
\usepackage[fontsize=10.5pt]{scrextend}
\usepackage[margin=22mm,headheight=15pt,headsep=9mm,footskip=11mm]{geometry}
\usepackage{fontspec}
\setmainfont{texgyrepagella}[Extension=.otf,UprightFont=*-regular,BoldFont=*-bold,ItalicFont=*-italic,BoldItalicFont=*-bolditalic]
\setsansfont{texgyreheros}[Extension=.otf,UprightFont=*-regular,BoldFont=*-bold,ItalicFont=*-italic,BoldItalicFont=*-bolditalic]
\setmonofont{texgyrecursor}[Extension=.otf,UprightFont=*-regular,BoldFont=*-bold,ItalicFont=*-italic,BoldItalicFont=*-bolditalic,Scale=MatchLowercase]
\usepackage{amsmath,amssymb}
\usepackage{unicode-math}
\setmathfont{texgyrepagella-math.otf}
\usepackage{xcolor,microtype,enumitem,needspace,fancyhdr}
\usepackage{bookmark}
\definecolor{ink}{HTML}{17344B}
\definecolor{accent}{HTML}{197D80}
\definecolor{muted}{HTML}{596773}
\hypersetup{colorlinks=true,linkcolor=accent,urlcolor=accent,pdftitle={MI-21 - Freewan--Hayajneh
inequality for sums of weighted geometric
means},pdfauthor={Open Problems in Numerical Linear Algebra}}
\urlstyle{same}
\setlength{\parindent}{0pt}
\setlength{\parskip}{5pt plus 1pt minus 1pt}
\linespread{1.04}
\setlength{\emergencystretch}{3em}
\widowpenalty=10000
\clubpenalty=10000
\displaywidowpenalty=10000
\allowdisplaybreaks[1]
\setlist{nosep,leftmargin=1.5em}
\providecommand{\tightlist}{\setlength{\itemsep}{0pt}\setlength{\parskip}{0pt}}
\providecommand{\pandocbounded}[1]{#1}
\pagestyle{fancy}
\fancyhf{}
\fancyhead[L]{\sffamily\footnotesize\color{muted}OPEN PROBLEMS IN NUMERICAL LINEAR ALGEBRA}
\fancyhead[R]{\sffamily\bfseries\footnotesize\color{accent}MI-21}
\fancyfoot[L]{\parbox[b]{0.9\textwidth}{\sffamily\fontsize{8}{10}\selectfont\color{muted}Editorial ratings; a bounded search cannot certify that a problem remains unsolved.\\Literature check: 2026-09-10}}
\fancyfoot[R]{\sffamily\footnotesize\color{muted}\thepage}
\renewcommand{\headrulewidth}{0.3pt}
\renewcommand{\headrule}{\hbox to\headwidth{\color{accent}\leaders\hrule height \headrulewidth\hfill}}
\makeatletter
\renewcommand{\section}{\Needspace{5\baselineskip}\@startsection{section}{1}{0pt}{12pt plus 2pt minus 2pt}{4pt}{\sffamily\bfseries\large\color{ink}}}
\renewcommand{\subsection}{\Needspace{4\baselineskip}\@startsection{subsection}{2}{0pt}{9pt plus 1pt minus 1pt}{3pt}{\sffamily\bfseries\normalsize\color{ink}}}
\makeatother
\setcounter{secnumdepth}{0}
\begin{document}
{\sffamily\small\color{muted}Matrix inequalities and norms\par}
\vspace{3pt}
{\raggedright\hyphenpenalty=10000\exhyphenpenalty=10000\sffamily\bfseries\fontsize{21}{24}\selectfont\color{ink}Freewan--Hayajneh
inequality for sums of weighted geometric means\par}
\vspace{7pt}
{\sffamily\small\textbf{Difficulty:} challenging\quad\textbf{Importance:} interesting
to the community\par}
{\sffamily\small\bfseries\color{accent}Status: Partially resolved\par}
\vspace{4pt}
{\color{accent}\hrule height 0.8pt}
\vspace{6pt}
\textbf{Rating rationale:} The coupled power and weight parameters make
extension beyond the proved cases challenging; positive definite
aggregation and matrix-function estimates give community relevance.

\subsection{Problem statement}\label{problem-statement}

For positive definite matrices, define \[
A\#_tB=A^{1/2}(A^{-1/2}BA^{-1/2})^tA^{1/2}.
\] Is the following true for all positive integers \(m,n\), all
\(A_i,B_i\in\mathbb C^{n\times n}\) positive definite, all
\(t\in[0,1]\), and all real \(s,r,p>0\) with \(sr\ge1\)? Set
\(A=\sum_iA_i\), \(B=\sum_iB_i\). For every unitarily invariant norm on
\(n\times n\) matrices, does \[
\left\|\sum_{i=1}^m(A_i^s\#_tB_i^s)^r\right\|
\le
\left\|\left(A^{(1-t)srp/2}B^{tsrp}A^{(1-t)srp/2}\right)^{1/p}\right\|
\] hold? Here all real powers of positive definite matrices use spectral
functional calculus, and unitary invariance means \(\|UXV\|=\|X\|\) for
unitary \(U,V\).

The conjecture controls the norm of a sum of nonlinear positive definite
matrix means using only sums of the input matrices. It is relevant to
aggregation and interpolation of positive definite data and to estimates
for algorithms involving matrix powers. The full weighted statement is
one problem; its individual parameter cases are not separate entries.

\subsection{References}\label{references}

\begin{enumerate}
\def\labelenumi{\arabic{enumi}.}
\tightlist
\item
  S. Freewan and M. Hayajneh, \emph{On a conjecture related to the
  geometric mean and norm inequalities}, Math. Inequal. Appl. 27 (2024),
  Conjecture 3, p.~195.
  \href{https://files.ele-math.com/articles/mia-27-15.pdf}{Publisher
  PDF}.
\item
  S. Freewan and M. Hayajneh, \emph{Norm inequalities involving
  geometric means}, arXiv:2401.00337v1, Conjecture 1.4 and the main
  theorem. \href{https://arxiv.org/abs/2401.00337}{Preprint}.
\end{enumerate}

Status check (2026-09-10): the current arXiv record lists v1, December
30, 2023. Its theorem treats \(t=1/2\), \(s\ge2\), \(r\ge1\), \(p>0\),
\(rp\ge1\); the abstract explicitly identifies this as a partial answer
to the weighted conjecture. The original publisher text and the later
preprint statement were compared. Searches for both exact titles, the
authors and ``weighted geometric mean conjecture'', including 2025/2026,
found no full proof or counterexample. Openness is subject to this
search limit.
\end{document}
